% ============================================================
%  Book preamble — Modern design with Inter + Fira Code
%  Included via --include-in-header (appended to pandoc's preamble)
% ============================================================

% ── Additional heading font families (optical-size variants) ─────
% pandoc sets \setmainfont{Inter 18pt} from book.yaml.
% We declare larger optical sizes for chapter/section headings.
\newfontfamily\headingfontlg{Inter 28pt}[
  BoldFont      = * Bold,
  ItalicFont    = * Italic,
  BoldItalicFont= * Bold Italic,
]
\newfontfamily\headingfontmd{Inter 24pt}[
  BoldFont      = * Bold,
  ItalicFont    = * Italic,
  BoldItalicFont= * Bold Italic,
]

% ── Accent color (used in headings and callout borders) ──────────
\usepackage{xcolor}
\definecolor{accent}{HTML}{ac4142}    % Base16 keyword red
\definecolor{accentblue}{HTML}{6a9fb5}% Base16 type blue
\definecolor{accentgreen}{HTML}{90a959}% Base16 variable green
\definecolor{accentorange}{HTML}{d28445}% Base16 definition orange
\definecolor{rulecolor}{HTML}{e2e8f0} % Light gray for rules
\definecolor{textmuted}{HTML}{718096} % Muted text

% ── Typography ───────────────────────────────────────────────────
\usepackage{microtype}
\usepackage{booktabs}
\usepackage{caption}
\captionsetup{font=small, labelfont={bf,sf}, labelsep=period}

% ── Modern table styling ─────────────────────────────────────────
\usepackage{array}        % better column specs + cell padding
\usepackage{colortbl}     % row/cell background colors

% More breathing room in cells
\renewcommand{\arraystretch}{1.45}
\setlength{\tabcolsep}{10pt}

% Muted rule color to match the overall palette
\arrayrulecolor{rulecolor}

% Heavier top/bottom rules from booktabs — keep them visible but subtle
\setlength{\heavyrulewidth}{0.8pt}
\setlength{\lightrulewidth}{0.4pt}
\setlength{\cmidrulewidth}{0.4pt}

% Header row color — applied via \rowcolor in a Lua filter;
% defined here so it's available as a named color throughout.
\definecolor{tableheadbg}{HTML}{eef2f7}   % cool light blue-gray
\definecolor{tableheadfg}{HTML}{2d3748}   % near-black for header text
\definecolor{tablerowalt}{HTML}{f8fafc}   % very subtle alternating row tint

% Redefine longtable header style injected by the Lua table filter.
% The filter wraps header cells in \thead{} — defined here:
\newcommand{\thead}[1]{%
  \textcolor{tableheadfg}{\bfseries\sffamily\small #1}%
}

% ── Code listing style ───────────────────────────────────────────
% ============================================================
%  Code listing style — fvextra line numbers + codeOS window frame
%
%  Included from preamble.tex via \input.
%  Depends on: xcolor (accent colors), tcolorbox [skins,breakable]
%  — both loaded earlier in preamble.tex.
% ============================================================

% ── Line numbers + soft wrap (fvextra) ───────────────────────────
% fvextra patches fancyvrb to support breaklines and line numbers.
% Must be loaded after fancyvrb (which pandoc's template loads first).
\usepackage{fvextra}

% Line number style: small, muted gray
\renewcommand{\theFancyVerbLine}{%
  \small\textcolor{textmuted}{\arabic{FancyVerbLine}}%
}

% Global defaults for all fancyvrb Verbatim environments
% (including pandoc's Highlighting environment).
\fvset{
  breaklines    = true,
  breakanywhere = false,
  numbers       = left,
  numbersep     = 12pt,
  xleftmargin   = 2em,
}

% ── codeOS window frame ───────────────────────────────────────────
% pandoc wraps syntax-highlighted blocks in a Shaded environment.
% We redefine it as a tcolorbox with a title-bar containing three
% traffic-light dots (left) and an optional centered filename.

% Colors
\definecolor{shadecolor}{HTML}{ffffff}    % fallback for framed/snugshade
\definecolor{codeosred}{HTML}{FF5F57}      % traffic-light red
\definecolor{codeosyellow}{HTML}{FEBC2E}   % traffic-light yellow
\definecolor{codeosgreen}{HTML}{28C840}    % traffic-light green
\definecolor{codeoswinframe}{HTML}{DEDEDE} % border + title separator
\definecolor{codewinbg}{HTML}{FFFFFF}    % window background

% Filename command — injected by filters/preamble.lua before each code
% block that carries a filename= attribute; cleared after each block.
\newcommand{\currentcodefilename}{}
\newcommand{\codeblockfilename}[1]{\gdef\currentcodefilename{#1}}

% Invisible right-side balancer so the filename stays centred in the bar.
\newcommand{\codeosphantom}{%
  \phantom{\tikz[baseline=-0.6ex, inner sep=0pt]{%
    \fill (0,0) circle[radius=5pt];%
    \fill (16pt,0) circle[radius=5pt];%
    \fill (32pt,0) circle[radius=5pt];%
  }}%
}

\renewenvironment{Shaded}{%
  \begin{tcolorbox}[%
    enhanced,
    breakable,
    colback      = codewinbg,
    colframe     = codeoswinframe,
    colbacktitle = codewinbg,
    coltitle     = codewinbg,
    boxrule      = 0.8pt,
    titlerule    = 0.5pt,
    arc          = 10pt,
    toptitle     = 8pt,
    bottomtitle  = 8pt,
    top          = 6pt,
    bottom       = 6pt,
    left         = 6pt,
    right        = 6pt,
    title        = {%
      % Traffic-light dots — zero-width anchor so \hfill centres the filename
      \makebox[0pt][l]{%
        \tikz[baseline=-0.6ex, inner sep=0pt]{%
          \fill[codeosred]    (0,0)    circle[radius=5pt];%
          \fill[codeosyellow] (16pt,0) circle[radius=5pt];%
          \fill[codeosgreen]  (32pt,0) circle[radius=5pt];%
        }%
      }%
      % Filename centred via \hfill + phantom mirror on the right
      \hfill
      \ifx\currentcodefilename\empty\else
        {\ttfamily\footnotesize\color{textmuted}\currentcodefilename}%
      \fi
      \hfill\codeosphantom
    },
  ]%
}{%
  \end{tcolorbox}%
  \gdef\currentcodefilename{}%
}


% ── Blank page after title ───────────────────────────────────────
% With twoside, the title lands on page 1 (odd/recto). Patching
% \maketitle to call \cleardoublepage forces page 2 to be blank and
% the TOC to start on page 3 (next odd/recto).
\let\origmaketitle\maketitle
\renewcommand{\maketitle}{\origmaketitle\cleardoublepage}

% ── Page headers and footers (fancyhdr) ──────────────────────────
% Two-sided layout: headers and footers hug the outer edge.
%   Even pages (verso / left):  content on the LEFT  (outer edge)
%   Odd  pages (recto / right): content on the RIGHT (outer edge)
\usepackage{fancyhdr}

\pagestyle{fancy}
\fancyhf{}
% Header: chapter title on the outer edge of each page
\fancyhead[LE]{\small\sffamily\color{textmuted}\nouppercase{\leftmark}}
\fancyhead[RO]{\small\sffamily\color{textmuted}\nouppercase{\leftmark}}
% Footer: page number on the outer edge of each page
\fancyfoot[LE]{\small\sffamily\color{textmuted}\thepage}
\fancyfoot[RO]{\small\sffamily\color{textmuted}\thepage}

\renewcommand{\headrulewidth}{0.4pt}
\renewcommand{\headrule}{%
  {\color{rulecolor}\hrule width\headwidth height\headrulewidth}%
}
\renewcommand{\footrulewidth}{0pt}

% Plain style (chapter opening pages): no header, page number on outer edge
\fancypagestyle{plain}{%
  \fancyhf{}%
  \fancyfoot[LE,RO]{\small\sffamily\color{textmuted}\thepage}%
  \renewcommand{\headrulewidth}{0pt}%
}

% ── Chapter & section heading styles (titlesec) ──────────────────
\usepackage{titlesec}

% Chapter: large number + title, accent rule below
\titleformat{\chapter}[display]
  {\normalfont\headingfontlg}
  {\large\color{accent}\sffamily\MakeUppercase{\chaptername}\ \thechapter}
  {6pt}
  {\Huge\bfseries}
  [\vspace{6pt}{\color{rulecolor}\titlerule[1.5pt]}]

\titlespacing*{\chapter}{0pt}{0pt}{28pt}

% Section: Inter 24pt SemiBold with subtle left accent bar
\titleformat{\section}
  {\headingfontmd\Large\bfseries}
  {\thesection}
  {0.75em}
  {}

\titlespacing*{\section}{0pt}{24pt}{6pt}

% Subsection: Inter 18pt SemiBold, no number decoration
\titleformat{\subsection}
  {\normalfont\large\bfseries\sffamily}
  {\thesubsection}
  {0.65em}
  {}

\titlespacing*{\subsection}{0pt}{16pt}{4pt}

% ── Callout boxes (tcolorbox) ────────────────────────────────────
\usepackage[skins,breakable]{tcolorbox}
\usepackage{fontawesome5}

% Base style shared by all callout variants
\tcbset{
  callout/.style={
    enhanced jigsaw,
    breakable,
    left        = 10pt,
    right       = 8pt,
    top         = 6pt,
    bottom      = 6pt,
    boxrule     = 0pt,
    leftrule    = 3pt,
    arc         = 2pt,
    fonttitle   = \bfseries\sffamily\small,
    coltitle    = #1,
    colback     = #1!8!white,
    colframe    = #1,
    attach boxed title to top left = {
      yshift    = -1mm,
      xshift    = 6pt,
    },
    boxed title style = {
      colback   = #1!8!white,
      colframe  = #1!8!white,
      arc       = 2pt,
      boxrule   = 0pt,
    },
  }
}

% Callout box labels — override any of these in book.yaml via:
%   header-includes:
%     - \newcommand{\callouttip}{...}
% \providecommand is a no-op when the command is already defined,
% so header-includes definitions always win.
\providecommand{\callouttip}{\faLightbulb\enspace Tip}
\providecommand{\calloutwarning}{\faExclamationTriangle\enspace Warning}
\providecommand{\calloutdanger}{\faFire\enspace Danger}

% \begin{tipbox}[Optional title override]
\newtcolorbox{tipbox}[1][\callouttip]{
  callout=accentblue,
  title=#1,
}

% \begin{warningbox}[Optional title override]
\newtcolorbox{warningbox}[1][\calloutwarning]{
  callout=accentorange,
  title=#1,
}

% \begin{dangerbox}[Optional title override]
\newtcolorbox{dangerbox}[1][\calloutdanger]{
  callout=accent,
  title=#1,
}

% ── List spacing ─────────────────────────────────────────────────
% tighter lists for a professional look
\usepackage{enumitem}
\setlist{noitemsep, topsep=4pt, partopsep=0pt}
\setlist[itemize,1]{label=\textcolor{accent}{\textbullet}}
